\section{Antes de empezar}
\subsection*{Cómo entras}
\begin{enumerate}
  \item Entra a Traza FI con tu correo institucional de alumno de posgrado.
  \item En el menú abre \emph{Titulación} → \emph{Posgrado: obtener mi grado}.
  \item Antes de comenzar puedes abrir la página de cada modalidad y descargar su manual; la modalidad la eliges al capturar tus datos.
\end{enumerate}
\begin{tencuenta}
\begin{itemize}
  \item La fecha del examen debe quedar al menos a 10 días hábiles; se sugieren 15.
  \item La firma electrónica UAEMéx (la del SICDE) de cada sinodal debe estar vigente al menos un mes después del examen; la plataforma avisa si alguna no alcanza.
  \item Las fotografías, las tres memorias USB y el resumen impreso se entregan en persona el día del cotejo.
  \item Solo puede haber un trámite activo. Si inicias otro, antes se te explica qué pasará con el que tienes.
  \item Aprobado el examen, tienes 48 horas para registrar tu grado en el SICDE y capturar el folio.
\end{itemize}
\end{tencuenta}
\begin{algosalemal}
\begin{itemize}
  \item \textbf{Te devolvieron un documento.} Tu siguiente paso te lleva a él; lee la observación, corrige el PDF y pulsa \boton{Reemplazar}.
  \item \textbf{Tu director o un sinodal no encuentra el correo.} En la lista de firmas puedes reenviar el enlace o corregir su correo.
  \item \textbf{Una firma UAEMéx no alcanza.} Pide a esa persona que la renueve en \texttt{https://fise.uaemex.mx/} y actualice la fecha desde su enlace.
  \item \textbf{Elegiste mal o quieres dejar el trámite.} Al final de tu trámite abre \emph{¿Elegiste otra modalidad o quieres dejar este trámite?}: si nadie más intervino se borra; si no, queda en pausa en \emph{Abandonadas} y lo retomas cuando quieras.
  \item \textbf{Tu examen se aplazó.} El segundo momento se reabre: propones una nueva fecha y tu director y tu sínodo vuelven a firmar.
\end{itemize}
\end{algosalemal}
